\chapter{技術者倫理・意思決定}

\section{背景：技術的に可能なことを、実施すべきことへ結び付ける}
技術は人の目的を実現する手段である。しかし、実現できること、依頼者が望むこと、社会が受け入れられることは必ずしも一致しない。材料の変更で納期と原価を改善できても、強度の低下によって使用者へ危険を移してよいとはいえない。専門家の判断は、その技術を十分に理解せず、個別に同意する機会もない人へ影響する。この知識と力の非対称性が技術者の社会的責任の背景にある。

物理モデルは「この条件なら何が起こるか」を説明する。倫理的判断は、その結果を誰が受け、どの負担を許容し、どの手続で決めるかも扱う。事実の把握が不十分なら判断は誤るが、事実だけから価値の優先順位が自動的に出るわけでもない。技術的な計算と価値判断の双方を明示する必要がある。

\section{公衆の優先と専門職への信頼}
\label{theory:ethics:public}
公衆には、顧客だけでなく、施設の周辺住民、製品の使用者、将来世代なども含まれる。依頼者が了承しても第三者の安全を損なう行為が許されるとは限らない。専門職への信頼は、力量、事実に基づく説明、公正な履行を通じて成立する。

\paragraph{出題当時と現在の資料}
2011年改訂の技術士倫理綱領は、公衆の利益、真実性、誠実な履行、正当理由のない秘密の開示の禁止などを示す\parencite{ipej_ethics2011}。2023年版では基本事項に具体的な指針が付され、リスクの評価、代替案、業務内容の説明等を明示する\parencite{ipej_ethics2023}。過去問の語句を現在の番号へ無断で置換しない。

安全を最優先にするとは、「ゼロリスク」を無根拠に断言することでも、全ての新技術を一律に禁止することでもない。危険を調査し、必要な仕様と対策を満たす案を見いだし、残る不確実さを関係者へ説明することである。

\section{倫理原理：結果・義務・徳という異なる問い}
\label{theory:ethics:principles}
\subsection{結果を検討する視点}
功利主義は、行為によって生じる幸福や利益、害の結果を重視する。技術の便益だけでなく、事故、環境負荷、負担の移転まで評価対象を広げる。ただし、全体の利益を合計すると少数者の重い負担が隠れる可能性がある。結果の集計と分配の公平性は別に検討する必要がある。

\subsection{義務と人格を尊重する視点}
義務を重視する立場は、有益な結果が予想されても守るべき約束や権利があると考える。虚偽の報告によって仕事を守ることは、望む結果を得るために相手を単なる手段として扱う問題を含む。人の同意、尊厳、真実へのアクセスを結果の総量だけで処分しない。

\subsection{徳と専門職の人格を問う視点}
徳を重視する立場は、その選択を繰り返すとどのような技術者になるかを問う。誠実さ、慎重さ、勇気、他者への配慮は、規程の空欄を埋める日常的な判断にも働く。能力が高いことと、責任ある判断をする人格は同一ではない。

各視点は問いが異なるので、同じ事例について異なる懸念を発見する。大学の倫理教育資料も複数の視点を用いた検討を提示する\parencite{scu_ethics_framework}。一つの原理が必ず唯一の正解を与えると仮定せず、衝突する理由を明らかにする。

\section{意思決定を手続として組み立てる}
\label{theory:ethics:procedure}
個人の直観だけで結論を出すと、知らない条件や他者への影響が抜けやすい。問題を明確にし、確かめた事実と不明事項を分け、影響を受ける人々を洗い出し、複数の行動案を比較し、相談・実行・振り返りへ進める。このような検討の枠組みが倫理的判断を支える\parencite{davis_seven_steps}。

工程変更なら、まず安全と要求品質が満たされるかを検討する。その制約を満たす案について納期と原価を比較する。「安い」「早い」を最初に固定して、残る要求を都合よく変更しない。ただし、個別設問で選んだ納期と原価の順位を全業務の不変な法則にしない。

予算を扱うときは、成果だけでなく承認と記録も重要である。超過や遅延の見通しを共有し、変更を承認する権限を持つ人へ説明する。私費を投入しても所有・管理・会計の手続を省けるわけではない。使い切るために不要な支出を作らず、未消化額を適切に扱う。

数量化できない評価項目をゼロにしてよいともいえない。安全、組織の信用、社会的な受容、復旧の重要性などは、尺度の定義と評価理由を示して比較する。計数の容易さは重要性の大小ではない。

\section{エシックス・テスト：判断に別の視点を入れる}
\label{theory:ethics:tests}
テストは正答を自動的に返す算法ではなく、見落としを発見する問いである。
\begin{description}
 \item[普遍化可能性] 同じ行為を誰もが行うと、その行為が前提とする約束や社会制度は成り立つかを考える。
 \item[可逆性] 影響を受ける相手の立場に自分を置いても同じ判断を支持するかを考える。
 \item[美徳] この選択を繰り返すことが自分の専門職としての人格をどう形作るかを考える。
 \item[世評] 判断とその理由が公に知られた場合に説明できるかを考える。
 \item[専門家] 所属する専門職の規範や他の専門家の吟味に照らして評価する。
\end{description}
立場交換、人格、公開、専門職の評価という視点はDavisの意思決定教材でも区別される\parencite{davis_seven_steps}。世評を単なる人気投票にせず、美徳を一回の印象だけで判断しない。一つのテストを通過しても、他の責任が免除されるわけではない。

\section{判断を阻害する要因と道徳的自律}
\label{theory:ethics:obstacles}
無知は判断に必要な情報を欠くことであり、知らないことを自覚して調べる姿勢とは異なる。自己中心的志向は自分の立場だけで問題を切り取ることである。微視的視野は細部にとらわれ、周囲や長期の影響を見失うことをいう。細かい調査そのものは悪いことではない。

自己欺瞞は、都合の悪い事実を直視せず、他人のため、今回だけ、皆もしているなどの理由で自分の行為を正当化することである。責任を果たす結果を恐れると、問題を報告せずに済ませたい誘因が生じる。利他性や勇気を一時的な感情で終わらせず、理由と証拠に基づく行動へ結び付ける。

道徳的自律とは、判断の理由を自分で吟味し、その判断への責任を持つことである。専門家への相談や組織での協働と矛盾しない。問題になるのは、指示者の権威だけを理由に判断を放棄すること、集団で合意したからという理由で批判的な検討を省くことである。異論を聞き、事実と評価を点検できる環境が重要となる。

\section{情報開示：警告と秘密保持を具体的に調整する}
\label{theory:ethics:disclosure}
危険の可能性を知った場合、何を確認できているか、被害がどれほど迫るか、誰へどの情報を伝えれば危険を減らせるかを検討する。必要な警告や緊急措置を遅らせる理由として秘密保持を使ってはならない。一方で、「危険の可能性があるなら即時・無条件に全情報を公開する」とするのも、誤情報、個人情報、関係のない秘密を広める危険がある。

報告先と開示範囲、確認状況、緊急性を記録して判断する。内部の検討が有効に機能するか、外部への通報が必要かも条件による。過去問で手続不足の告発が不適切とされたことから、告発全般を否定する結論を導かない。法的な保護要件と具体的な通報手続は法令・説明責任の章で別に扱う。

\section{個人から組織へ：倫理規程と学習の仕組み}
\label{theory:ethics:organization}
不正防止を個人の心掛けだけに任せると、納期や売上の圧力、相談の難しさ、記録の不備が残る。記録と報告の経路、専門能力を持つ人の関与、異論を伝えられる環境、類似事例の共有を整える必要がある。対策会議を開催しても、実際の検討を現場へ押し付けるなら組織としての役割を果たさない。

倫理規程は社会への約束、目指す理想、行動の指針、教育、組織自体を点検する対話の基準になる。本文が存在するだけでは実効性を保証せず、具体的な状況に適用し、運用を見直すことが必要である。JSTの教育資料でも倫理綱領と継続的な組織の取組を関連付ける\parencite{fudano_ethics2021}。ここでの社会との約束は、個別の法律上の契約と同一ではない。

データの真正性を守ることは組織の信用の前提である。誤記の訂正や科学的に妥当な処理には根拠と記録が必要で、依頼者が望む結果に合わせた変更とは違う。顧客の要求や上司の指示は、事実を改変して報告する免許ではない。研究固有の記録と不正行為は研究倫理の章も参照する。

\section{予防倫理と志向倫理}
\label{theory:ethics:aspiration}
予防倫理は有害な行為や不正を防ぎ、守るべき境界を示す。志向倫理は、よりよい意思決定や社会への積極的な貢献へ向かう側面である。違反していないことだけでは、技術者として果たし得る役割を全て果たしたとは限らない。両方を認識して行動することが重要である\parencite{fudano_ethics2021}。

\paragraph{出題当時の公務員倫理}
2017年の人事院教材は、倫理法・規程の範囲を狭義、公務員としての在り方を含む範囲を広義として区別する\parencite{jinji_ethics2017}。平成30年度はこの範囲と技術者倫理の二つの側面を対応させて問う。利害関係者は職務との利害関係に基づく対象であり、単に知人という分類ではない。

\paragraph{現在の扱い}
人事院の説明は、利害関係者との贈与・接待やゴルフ・旅行等の規制を解説し、最新の事例集への参照も案内する\parencite{jinji_ethics_qa}。この教材の語句要約を、例外も含む全規制の説明に置き換えない。出題後のJST資料は予防・志向の理解を補う背景資料として使い、平成30年時点の資料だとは扱わない。

\section{社会的責任と共有価値}
\label{theory:ethics:csr}
企業は経済的な成果を生むとともに、従業員、利用者、地域、環境へ影響する。社会的責任の取組を、利益が全く生じない慈善活動だけに限定できない。有償の回収でも廃棄物削減へ貢献し得るし、社会貢献に伴う信用の向上が売上へ寄与することもあり得る。利益を得た事実だけで社会的効果が消えるわけではない。

共有価値の創出は、社会課題の解決と企業の経済的価値を事業の仕組みとして結び付ける構想である。提唱者所属の研究機関はCSRや寄付との概念差も明示する\parencite{hbs_csv}。CSRとCSVを同一の用語にせず、全ての寄付がCSVになるともいわない。過去問の候補では、両者の定義の細部より、利益と社会的責任が両立しないとの誤った限定が問われる場合がある。

\subsection{社会的責任のガイダンスと認証}
要求事項への適合を認証する仕組みと、組織の行動について指針を与える文書を区別する。
ISO26000は社会的責任のガイダンスであり、認証のための規格ではない\parencite{iso26000}。
H23の例では3を公式正答として保持するが、他肢の対象組織や原則列挙の表現には疑義が残る。当時の文言、公開問題、現在確認できる一次資料の範囲を分ける。一つの正答理由を確認しただけで全肢の検証を完了したとは扱わない。

\section{行動規範：専門能力の限界と未来世代への責任}
\label{theory:ethics:code}
電気学会の2007年行動規範は、技術者の能力には限界があるという認識から、多様な見識を持つ人々との協同を重視する。目標は技術進歩それ自体だけでなく、持続可能な社会と未来世代の生存、安全・健康・福祉への貢献である\parencite{ieej_conduct2007}。同じ専門分野・同じ価値観の人だけで判断すると、影響を受ける他者の事情や、技術外の制度・文化の条件が抜ける。

行動規範は具体的な行動の指針、倫理綱領は基本的な約束として読む。平成30年の前文問題を2007年規範へ対応させ、2021年改訂の綱領と同じ時点の文書とは扱わない。可逆性の点検でも、自分が負担を許容できたことを他者全員の同意へ置き換えず、異なる見識との対話を続ける。

\section{ISO26000：組織の影響を社会的責任へ結び付ける}
\label{theory:ethics:iso}
社会的責任は、組織の活動が社会や環境へ及ぼす影響を判断と行動に反映することである。組織は契約相手だけでなく、雇用者、利用者、地域などの利害関係者と関係する。利害関係者は株主と同義ではなく、非営利組織や行政機関にも存在する。ISO26000の対象は広く組織であり、著名人や公人を独立の適用対象として並列に加える定義ではない。個人が組織の活動に担う役割とは分ける。

2010年発行時のISO資料も、ガイダンスであって認証用の要求事項ではないこと、持続可能な発展への貢献を目的とすることを示す\parencite{iso26000_project2010}。対策を実施した結果を検討するために使用する文書と、適合証明のための要求事項を混同しない。

\subsection{行動の原則と扱う主題を分ける}
七原則は、説明責任、透明性、倫理的行動、利害関係者の利害の尊重、法の支配の尊重、国際行動規範の尊重、人権の尊重である。七つの中核主題は、組織統治、人権、労働慣行、環境、公正な事業慣行、消費者課題、コミュニティへの参画と発展であり、同じリストではない。公式の対応資料はこの二層を分けて示す\parencite{iso_gri2014}。

原則は判断の仕方、主題は検討対象を整理する。たとえば労働慣行について従業員の利害を把握し、判断を透明にし、結果への説明責任を持つというように組み合わせる。「法の支配」を権力者の好みに従うことへ置換すると、法によって権力を制約する意味が失われる。政府の意向を独立の尊重原則として付加する肢は、この区別で吟味できる。透明性も、秘密や個人情報を一律に無条件公開することとは区別する。

\section{NPO：非営利と無収入・無報酬は違う}
\label{theory:ethics:npo}
企業の収益事業だけが技術者の社会参加ではない。技術を用いた災害支援や地域での協働などには、非営利組織としての活動もある。特定非営利活動法人は、法定の活動分野と、不特定多数の利益の増進という目的の条件を持つ。2012年4月施行の改正後は20分野となり、平成24年度問題はこの時点の制度を問う\parencite{cabinet_npo_overview}。

非営利とは利益を構成員へ分配する目的を持たないことである。活動を続ける収入や剰余、職員への給与がない意味ではない。主目的の特定非営利活動に支障がない範囲のその他の事業も可能である。利益分配、公益の事業目的、労働への対価は異なる判断軸である。役員として報酬を受ける人数は役員総数の3分の1以下という条件があり、職員の労働への給与はこの役員報酬とは分ける\parencite{cabinet_npo_officers}。

法定分野の存在から、営利企業を競争から保護するためにNPOを制限する制度だとは導けない。会社員であること自体も参加禁止の理由にはならず、就業規則、利益相反、勤務時間、秘密保持などの具体的条件を点検する。

\section{仕事と生活の調和：選択できる環境を整える}
\label{theory:ethics:worklife}
職場の責任と、育児・介護・地域生活などの責任は、人生の段階によって重なり方が変わる。仕事と生活の調和は、その多様性を前提に、働き続けながら選択を実現できる条件を整える考え方である。2007年の憲章と2010年の改定を背景に、平成27年度の問題を読む\parencite{cabinet_worklife}。

柔軟性には勤務時間、場所、休暇、仕事の分担など複数の軸がある。雇用契約を短期にして不安定さを増すこととは同一ではない。短期非正規への置換を必要条件とする記述は、生活を安定させて選択を支える目的との整合を欠く。採用や定着、就業継続の改善は期待効果であり、企業ごとの実現を無条件に保証しない。国・自治体、労使、個人が条件を改善する役割を分担する。

\section{ユニバーサルデザイン：利用者を設計の前提へ入れる}
\label{theory:ethics:universal}
利用方法を一種類の身体能力や言語、経験に合わせると、それ以外の人は利用しづらくなる。設計段階から多様性を扱うには、利用の公平さ、方法の柔軟性、直観的な操作、必要情報の認知、エラーの害の低減、身体的負担、接近と使用の空間という七つの視点を点検する。ノースカロライナ州立大学の1997年版が七原則の原資料であり、現在のウェブやPDFの掲載年と混同しない\parencite{ncsu_ud1997}。

たとえば操作の結果を視覚と音で知らせる案なら、情報を受け取る条件の多様性を検討できる。エラーへの寛容は、誤操作を起こしてよいという指示ではなく、起こり得る誤操作が重い害へ直結しないよう設計することである。全ての人へ常に完全な適合を保証する標語として扱わず、利用者の条件と安全・費用等を評価する。既存の障壁を除くバリアフリーも必要であり、その対策が常に浪費であるとの一般則を設問の前文から導かない。

\section{出題で確かめる条件}
「即時・無条件」「要求があれば」「全て」「だけ」などの限定を確認する。前半が適切でも、後半に不正を認める例外があれば記述全体は適切でない。数を問う場合は各記述を先に判断してから集計し、組合せでは表の列と選択肢番号を区別する。

ここで扱った各過去問は、原典の条件、独立した判断、公式正答との比較、全5肢の理由へ追跡できる。次の問題ブロックはAstra査読前の初稿である。章内の未執筆問題を代表例で置き換えず、全問の処理を継続する。
\section{告発の倫理：段階と証拠、改善への見込み}
\label{theory:ethics:whistle_conditions}
深刻な危害が見込まれるとき、内部で懸念を伝え有効な是正を求める。
上司が応答しないという一事だけで、他の内部経路も調べず直ちに外部公表へ進むとはしない。
客観的な証拠、危害を伝える相手、公表による改善の見込みと個人に及ぶ危険を考える。
これは倫理的な意思決定を整理する枠組みで、法定の公益通報保護要件そのものではない。\cite{muroran_ethics_cases}

令和元年度再試験の比較では、証拠がなくても外部へ知らせる義務が必ずあるという肢を否定する。
一方、危害が急迫している実際の事態に、全ての内部段階を必ず順番に待つという手順を命じるものでもない。
緊急性と有効な通報先、必要な開示範囲、現在の法定経路を調べる。
原著全章を確認した解説ではないため、倫理モデルの解釈は査読事項として保持する。

\section{宣言の三つの責任}
\label{theory:profession:declaration}
2007年のプロフェッション宣言は、自らの技術、顧客、社会に対する責任を示す。能力を高め、正当な守秘を保ちつつ顧客へ誠実に業務を行い、社会と環境への影響を考える\parencite{eth54_profession2007}。顧客だけ、所属企業だけへ責任を狭めない。

倫理綱領は専門職の行動を支える規範である。団体の会員上の適用と、法律・契約・組織規程の適用は別に考える。非会員だから別の義務がすべてなくなるわけではない。反対に、特定団体の内部規範をすべての他団体へそのまま強制できるとも仮定しない。

\section{多様性とDEI：能力を発揮できる条件を整える}
\label{theory:ethics:diversity}
多様性には性別や国籍などの属性だけでなく、知識、経験、専門性、価値観、働き方の違いもある。人の違いを受け入れ、参加の障壁を調べる。違いの尊重は、他者へ危害を与える行動まで無条件に認めることではない。

日本技術士会の2025年宣言は、多様性、公平性、包摂性の三つを結びつける\parencite{eth54_dei2025}。公平性では、一人ひとりの条件に応じて環境や道具を用意する。同じ道具を配るだけでは利用の障壁が残る場合がある。包摂性は、違いを理由に排除せず、能力を発揮して組織の価値を生む状態を目指す。2021年の多様性問題に用いる場合は、後年の補足資料であることを明示する。

\subsection{平成23年度 問11}
\label{q:h23-reg-apt-11}

\noindent\textbf{原典：} sources/H23/aptitude.pdf、PDF物理12頁。

\textbf{公式正答：} 3（single）。

\paragraph{問題の要旨} 社会的責任の手引きと認証規格の区別を問う。公式正答③を保持し、他の記述の解釈上の疑義を別に残す。

\paragraph{必要知識} 社会的責任、要求事項とガイダンス、認証の適用範囲

\paragraph{考え方（査読待ち初稿）} ISO26000は社会的責任に関するガイダンスであり、認証のための要求事項を規定した規格ではない。③がこの性格を誤っているため公式正答は③。ISOの一次資料でも認証用途ではないことを確認した。

\paragraph{他の選択肢（査読待ち初稿）} 他肢をすべて無条件に正しいと説明しない。①の政府除外の表現、④の原則の列挙範囲には一次資料との照合を要する疑義が残る。②⑤の歴史的な参加・連携の詳細も当時資料の確認を後工程に残す。

\par 参考：\cite{iso26000}。

\par 難易度 3（校準初稿）、優先度未評価。

\subsection{平成24年度 問2}
\label{q:h24-reg-apt-02}

\noindent\textbf{原典：} sources/H24/aptitude.pdf、PDF物理1頁。

\textbf{公式正答：} 3（single）。

\par\textit{原典本文の校正と独自解説は未作成。分類のみの登録項目。}

\subsection{平成24年度 問12}
\label{q:h24-reg-apt-12}

\noindent\textbf{原典：} sources/H24/aptitude.pdf、PDF物理10頁。

\textbf{公式正答：} 2（single）。

\par\textit{原典本文の校正と独自解説は未作成。分類のみの登録項目。}

\subsection{平成24年度 問13}
\label{q:h24-reg-apt-13}

\noindent\textbf{原典：} sources/H24/aptitude.pdf、PDF物理11頁。

\textbf{公式正答：} 3（single）。

\par\textit{原典本文の校正と独自解説は未作成。分類のみの登録項目。}

\subsection{平成25年度 問6}
\label{q:h25-reg-apt-06}

\noindent\textbf{原典：} sources/H25/aptitude.pdf、PDF物理7頁。

\textbf{公式正答：} 5（single）。

\par\textit{原典本文の校正と独自解説は未作成。分類のみの登録項目。}

\subsection{平成25年度 問12}
\label{q:h25-reg-apt-12}

\noindent\textbf{原典：} sources/H25/aptitude.pdf、PDF物理13頁。

\textbf{公式正答：} 5（single）。

\par\textit{原典本文の校正と独自解説は未作成。分類のみの登録項目。}

\subsection{平成26年度 問6}
\label{q:h26-reg-apt-06}

\noindent\textbf{原典：} sources/H26/aptitude.pdf、PDF物理6頁。

\textbf{公式正答：} 5（single）。

\par\textit{原典本文の校正と独自解説は未作成。分類のみの登録項目。}

\subsection{平成27年度 問3}
\label{q:h27-reg-apt-03}

\noindent\textbf{原典：} sources/H27/aptitude.pdf、PDF物理3頁。

\textbf{公式正答：} 2（single）。

\par\textit{原典本文の校正と独自解説は未作成。分類のみの登録項目。}

\subsection{平成27年度 問6}
\label{q:h27-reg-apt-06}

\noindent\textbf{原典：} sources/H27/aptitude.pdf、PDF物理6頁。

\textbf{公式正答：} 3（single）。

\par\textit{原典本文の校正と独自解説は未作成。分類のみの登録項目。}

\subsection{平成27年度 問15}
\label{q:h27-reg-apt-15}

\noindent\textbf{原典：} sources/H27/aptitude.pdf、PDF物理16頁。

\textbf{公式正答：} 4（single）。

\par\textit{原典本文の校正と独自解説は未作成。分類のみの登録項目。}

\subsection{平成29年度 問3}
\label{q:h29-reg-apt-03}

\noindent\textbf{原典：} sources/H29/aptitude.pdf、PDF物理3頁。

\textbf{公式正答：} 3（single）。

\par\textit{原典本文の校正と独自解説は未作成。分類のみの登録項目。}

\subsection{平成29年度 問7}
\label{q:h29-reg-apt-07}

\noindent\textbf{原典：} sources/H29/aptitude.pdf、PDF物理7頁。

\textbf{公式正答：} 4（single）。

\par\textit{原典本文の校正と独自解説は未作成。分類のみの登録項目。}

\subsection{平成29年度 問15}
\label{q:h29-reg-apt-15}

\noindent\textbf{原典：} sources/H29/aptitude.pdf、PDF物理14頁。

\textbf{公式正答：} 2（single）。

\par\textit{原典本文の校正と独自解説は未作成。分類のみの登録項目。}

\subsection{平成30年度 問4}
\label{q:h30-reg-apt-04}

\noindent\textbf{原典：} sources/H30/aptitude.pdf、PDF物理4頁。

\textbf{公式正答：} 1（single）。

\par\textit{原典本文の校正と独自解説は未作成。分類のみの登録項目。}

\subsection{平成30年度 問14}
\label{q:h30-reg-apt-14}

\noindent\textbf{原典：} sources/H30/aptitude.pdf、PDF物理14頁。

\textbf{公式正答：} ALL（all\_awarded）。

\par\textit{原典本文の校正と独自解説は未作成。分類のみの登録項目。}

\subsection{平成30年度 問15}
\label{q:h30-reg-apt-15}

\noindent\textbf{原典：} sources/H30/aptitude.pdf、PDF物理15頁。

\textbf{公式正答：} 4（single）。

\par\textit{原典本文の校正と独自解説は未作成。分類のみの登録項目。}

\subsection{令和元年度 問9}
\label{q:r01-reg-apt-09}

\noindent\textbf{原典：} sources/R01/aptitude.pdf、PDF物理10頁。

\textbf{公式正答：} 2（single）。

\par\textit{原典本文の校正と独自解説は未作成。分類のみの登録項目。}

\subsection{令和元年度 問14}
\label{q:r01-reg-apt-14}

\noindent\textbf{原典：} sources/R01/aptitude.pdf、PDF物理14頁。

\textbf{公式正答：} 4（single）。

\par\textit{原典本文の校正と独自解説は未作成。分類のみの登録項目。}

\subsection{令和元年度再試験 問3}
\label{q:r01-ret-apt-03}

\noindent\textbf{原典：} sources/R01\_retest/aptitude.pdf、PDF物理3頁。

\textbf{公式正答：} 1（single）。

\par\textit{原典本文の校正と独自解説は未作成。分類のみの登録項目。}

\subsection{令和2年度 問2}
\label{q:r02-reg-apt-02}

\noindent\textbf{原典：} sources/R02/aptitude.pdf、PDF物理2頁。

\textbf{公式正答：} 2（single）。

\par\textit{原典本文の校正と独自解説は未作成。分類のみの登録項目。}

\subsection{令和2年度 問11}
\label{q:r02-reg-apt-11}

\noindent\textbf{原典：} sources/R02/aptitude.pdf、PDF物理11頁。

\textbf{公式正答：} 1（single）。

\par\textit{原典本文の校正と独自解説は未作成。分類のみの登録項目。}

\subsection{令和2年度 問15}
\label{q:r02-reg-apt-15}

\noindent\textbf{原典：} sources/R02/aptitude.pdf、PDF物理15頁。

\textbf{公式正答：} 1（single）。

\par\textit{原典本文の校正と独自解説は未作成。分類のみの登録項目。}

\subsection{令和5年度 問2}
\label{q:r05-reg-apt-02}

\noindent\textbf{原典：} sources/R05/aptitude.pdf、PDF物理2頁。

\textbf{公式正答：} 3（single）。

\par\textit{原典本文の校正と独自解説は未作成。分類のみの登録項目。}

\subsection{令和5年度 問9}
\label{q:r05-reg-apt-09}

\noindent\textbf{原典：} sources/R05/aptitude.pdf、PDF物理9頁。

\textbf{公式正答：} 4（single）。

\par\textit{原典本文の校正と独自解説は未作成。分類のみの登録項目。}

\subsection{令和5年度 問11}
\label{q:r05-reg-apt-11}

\noindent\textbf{原典：} sources/R05/aptitude.pdf、PDF物理11頁。

\textbf{公式正答：} 2（single）。

\par\textit{原典本文の校正と独自解説は未作成。分類のみの登録項目。}

\subsection{令和7年度 問3}
\label{q:r07-reg-apt-03}

\noindent\textbf{原典：} sources/R07/aptitude.pdf、PDF物理3頁。

\textbf{公式正答：} 5（single）。

\par\textit{原典本文の校正と独自解説は未作成。分類のみの登録項目。}

